
\documentclass[12pt]{article}
\usepackage[utf8]{inputenc}
\usepackage[portuguese]{babel}
\usepackage{booktabs}
\usepackage{array}
\usepackage{geometry}
\usepackage{amsmath}
\usepackage{amsfonts}
\usepackage{longtable}

\geometry{a4paper, margin=1.5cm}

\title{Solução do Problema Afiro - Solver HiGHS}
\author{Análise Computacional}
\date{\today}

\begin{document}

\maketitle

\section{Informações do Problema}

\textbf{Informações do Problema:}
\begin{itemize}
\item Nome: Afiro
\item Número de Variáveis: 32
\item Número de Restrições: 27
\item Inviabilidade Primal: 0.000e+00
\item Inviabilidade Dual: 0.000e+00
\item Valor Primal: -4.648e+02
\item Valor Dual: -4.648e+02
\item Gap: -1.314e-19
\item Número de Iterações: 7
\end{itemize}


\section{Variáveis Primais e Custos Reduzidos (x > 0 e z > 0)}

\begin{longtable}{@{}cccc@{}}
\caption{Variáveis primais (x) e custos reduzidos (z) com valores > 0 do problema Afiro} \\
\toprule
\textbf{Coordenada x} & \textbf{Valor x (Primal)} & \textbf{Coordenada z} & \textbf{Valor z (Custo Reduzido)} \\
\midrule
\endfirsthead

\toprule
\textbf{Coordenada x} & \textbf{Valor x (Primal)} & \textbf{Coordenada z} & \textbf{Valor z (Custo Reduzido)} \\
\midrule
\endhead

\midrule \multicolumn{4}{r}{{Continua na próxima página}} \\ \midrule
\endfoot

\bottomrule
\endlastfoot
1 & 8.000e+01 & 6 & 2.250e+00 \\
2 & 2.550e+01 & 7 & 2.270e+00 \\
3 & 5.450e+01 & 8 & 2.290e+00 \\
4 & 8.480e+01 & 9 & 2.229e+00 \\
5 & 8.000e+01 & 25 & 2.066e+00 \\
13 & 1.821e+01 & 26 & 2.092e+00 \\
14 & 6.179e+01 & 27 & 2.120e+00 \\
15 & 8.480e+01 & 28 & 2.149e+00 \\
16 & 5.000e+02 & 32 & 1.000e+01 \\
17 & 4.759e+02 &  &  \\
18 & 2.408e+01 &  &  \\
20 & 2.150e+02 &  &  \\
29 & 3.399e+02 &  &  \\
30 & 3.839e+02 &  &  \\

\bottomrule
\end{longtable}


\section{Variáveis Duais (Multiplicadores de Lagrange - Valores > 0)}

\begin{longtable}{@{}cc@{}}
\caption{Variáveis duais (y) com valores > 0 do problema Afiro} \\
\toprule
\textbf{Coordenada y} & \textbf{Valor y (Dual)} \\
\midrule
\endfirsthead

\toprule
\textbf{Coordenada y} & \textbf{Valor y (Dual)} \\
\midrule
\endhead

\midrule \multicolumn{2}{r}{{Continua na próxima página}} \\ \midrule
\endfoot

\bottomrule
\endlastfoot

\bottomrule
\end{longtable}


\end{document}
